% Part: first-order-logic
% Chapter: completeness
% Section: construction-of-model

\documentclass[../../../include/open-logic-section]{subfiles}

\begin{document}

\iftag{FOL}
      {\olfileid{fol}{com}{mod}}
      {\olfileid{pl}{com}{mod}}

\olsection{একটি মডেলের নির্মাণ}

\begin{explain}
\iftag{FOL}{এখন আমরা~$\eq$ নিয়ে ভাবছি না; অর্থাৎ শুধু দেখাতে চাই যে
  $\eq$-বিহীন !!{sentence}s-এর কোনো সঙ্গত সেট~$\Gamma$ পরিতৃপ্তিযোগ্য।
  প্রথমে~$\Gamma$-কে একটি সঙ্গত, !!{complete} ও সম্পৃক্ত
  সেট~$\Gamma^*$-তে প্রসারিত করি। এক্ষেত্রে মডেল~$\Struct{M(\Gamma^*)}$-এর
  সংজ্ঞা সরল: $\Lang{L'}$-এর বদ্ধ পদগুলির সেটকে সংজ্ঞাক্ষেত্র নিই। প্রতিটি
  !!{constant}-কে নিজেকেই দিই, এবং আরও সাধারণভাবে নিশ্চিত করি যে প্রতিটি
  বদ্ধ পদ~$t$-এর জন্য $\Value{t}{M(\Gamma^*)} = t$।
  !!{predicate}s-কে এমনভাবে বিস্তার দিই যাতে কোনো পরমাণু !!{sentence},
  $\Struct{M(\Gamma^*)}$-তে সত্য হয় যদি এবং কেবল যদি সেটি~$\Gamma^*$-তে
  থাকে। এতে স্পষ্টতই~$\Gamma^*$-র সব পরমাণু !!{sentence}s,
  $\Struct{M(\Gamma^*)}$-তে সত্য হবে। যে~$\Gamma^*$ দিয়ে শুরু করি সেটি
  সঙ্গত, পূর্ণ ও সম্পৃক্ত হলে অবশিষ্ট বাক্যগুলিও সত্য হবে।}{এবার
  এমন একটি !!a{valuation} সংজ্ঞায়িত করার জন্য আমরা প্রস্তুত, যা প্রতিটি
  $!A \in \Gamma$-কে সত্য করে। এজন্য প্রথমে লিন্ডেনবাউমের সহায়ক উপপাদ্য
  প্রয়োগ করি: একটি পূর্ণ সঙ্গত $\Gamma^* \supseteq \Gamma$ পাই।
  $\Gamma^*$-র !!{propositional variable}s দিয়ে~$\pAssign
  v(\Gamma^*)$ নির্ধারণ করি।}
\end{explain}

\iftag{FOL}{%
\begin{defn}[পদ-মডেল]
\ollabel{defn:termmodel}
ধরা যাক কোনো ভাষা~$\Lang L$-এর !!{sentence}s-এর সেট $\Gamma^*$
!!a{complete}, সঙ্গত ও সম্পৃক্ত। $\Gamma^*$-র \emph{পদ-মডেল}~$\Struct
M(\Gamma^*)$ নিচের মতো সংজ্ঞায়িত !!{structure}:
\begin{enumerate}
\item !!{domain}~$\Domain{M(\Gamma^*)}$ হলো~$\Lang L$-এর সব বদ্ধ পদের
  সেট।
\item কোনো !!a{constant}~$c$-এর ব্যাখ্যা নিজেই $c$:
  $\Assign{c}{M(\Gamma^*)} = c$।
\item !!{function}~$f$-কে এমন অপেক্ষকটি দেওয়া হয়, যা যুক্তি হিসেবে
  বদ্ধ পদ $t_1$, \dots, $t_n$ পেলে বদ্ধ পদ $f(t_1, \dots, t_n)$-কে মান
  হিসেবে দেয়:
\[
\Assign{f}{M(\Gamma^*)}(t_1, \dots, t_n) = f(t_1,\dots, t_n)
\]
\item $R$ একটি $n$-স্থানীয় !!{predicate} হলে
  \[
  \tuple{t_1, \dots,
  t_n} \in \Assign{R}{M(\Gamma^*)} \text{ যদি এবং কেবল যদি } \Atom{R}{t_1, \dots,
    t_n} \in \Gamma^*.
  \]
\end{enumerate}
\end{defn}
}{%
\begin{defn}
\ollabel{defn:termmodel}
ধরা যাক $\Gamma^*$ হলো !!{formula}s-এর একটি পূর্ণ সঙ্গত সেট। তখন রাখি
\[
\pAssign v(\Gamma^*)(p) = \begin{cases}
  \True & \text{$p \in \Gamma^*$ হলে}\\
  \False & \text{$p \notin \Gamma^*$ হলে}
  \end{cases}
\]
\end{defn}
}

\iftag{FOL}{%
এবার যাচাই করব যে সত্যিই $\Value{t}{M(\Gamma^*)} = t$।

\begin{lem}
\ollabel{lem:val-in-termmodel} $\Struct M(\Gamma^*)$-কে
\olref{defn:termmodel}-এর পদ-মডেল ধরি। তাহলে প্রতিটি বদ্ধ পদের জন্য
$\Value{t}{M(\Gamma^*)} = t$।% BN-SRC-076 TARGET-CORRECTION-BEGIN
\par\smallskip\noindent\textnormal{\small\textbf{উৎস-স্পষ্টীকরণ (BN-SRC-076):} হিমায়িত ইংরেজি সহায়ক উপপাদ্যের সিদ্ধান্তে $t$-কে অবাধ রাখা হয়েছে, যদিও পদ-মডেলের এই মানায়নটি চল-মূল্যায়ন ছাড়া কেবল বদ্ধ পদের জন্য সংজ্ঞায়িত এবং প্রমাণের ভিত্তি ও আরোহী ধাপ উভয়েই বদ্ধ পদ বলা হয়েছে। বাংলা লক্ষ্যে সিদ্ধান্তেও বদ্ধ-পদ পরিসরটি স্পষ্ট করা হয়েছে।} % BN-SRC-076 TARGET-CORRECTION-END
\end{lem}
\begin{proof}
 $t$-এর ওপর আরোহের সাহায্যে প্রমাণ করি। $t$ একটি !!a{constant} হলে
 ভিত্তিটি পদ-মডেলের সংজ্ঞা থেকে সরাসরি অনুসৃত হয়। আরোহী ধাপের জন্য ধরা
 যাক $t_1, \ldots, t_n$ এমন বদ্ধ পদ যাতে
 $\Value{t_i}{M(\Gamma^*)} = t_i$, এবং $f$ একটি $n$-স্থানীয়
 !!{function}। তাহলে
\begin{align*}
\Value{f(t_1,\ldots,t_n)}{M(\Gamma^*)} &= \Assign{f}{M(\Gamma^*)}(\Value{t_1}
 {M(\Gamma^*)},
\ldots, \Value{t_n}{M(\Gamma^*)}) \\
&= \Assign{f}{M(\Gamma^*)}(t_1, \dots, t_n) \\
&= f(t_1,\dots, t_n),
\end{align*}
তাই আরোহ অনুসারে প্রতিটি বদ্ধ পদ~$t$-এর জন্য কথাটি সত্য।
\end{proof}
}

\iftag{FOL}{%
\begin{explain}
  \iftag{prvEx}{কোনো !!^a{structure}~$\Struct{M}$ একটি অস্তিত্বমূলকভাবে
    পরিমাণিত !!{sentence}~$\lexists[x][!A(x)]$-কে সত্য করতে পারে, যদিও
    সেটি সত্য করে এমন কোনো নিদর্শন~$!A(t)$ না-ও থাকতে পারে।}{}
  \iftag{prvAll}{কোনো !!^a{structure}~$\Struct{M}$ সার্বিকভাবে পরিমাণিত
    !!{sentence}~$\lforall[x][!A(x)]$-এর সব নিদর্শন~$!A(t)$-কে সত্য করতে
    পারে, কিন্তু~$\lforall[x][!A(x)]$-কে সত্য না-ও করতে পারে।}{} কারণ
  সাধারণভাবে~$\Domain{M}$-এর প্রতিটি !!{element} কোনো বদ্ধ পদের মান নয়
  ($\Struct{M}$ আচ্ছাদিত না-ও হতে পারে)। এই কারণেই পরিতৃপ্তি-সম্পর্ককে
  চল-মূল্যায়নের সাহায্যে সংজ্ঞায়িত করা হয়। কিন্তু আমাদের
  পদ-মডেল~$\Struct{M(\Gamma^*)}$-এর ক্ষেত্রে তার দরকার হতো না—কারণ সেটি
  আচ্ছাদিত। পরের ফলটি ঠিক এই কথাই বলে।
\end{explain}

\begin{prop}
\ollabel{prop:quant-termmodel}
$\Struct M(\Gamma^*)$-কে \olref{defn:termmodel}-এর পদ-মডেল ধরি।
\begin{tagenumerate}{prvEx,prvAll}
\tagitem{prvEx}{$\Sat{M(\Gamma^*)}{\lexists[x][!A(x)]}$ যদি এবং কেবল
  যদি অন্তত একটি বদ্ধ পদ~$t$-এর জন্য $\Sat{M(\Gamma^*)}{!A(t)}$।}{}
\tagitem{prvAll}{$\Sat{M(\Gamma^*)}{\lforall[x][!A(x)]}$ যদি এবং কেবল
  যদি সব বদ্ধ পদ~$t$-এর জন্য $\Sat{M(\Gamma^*)}{!A(t)}$।}{}
\end{tagenumerate}
\end{prop}

\begin{proof}
\begin{tagenumerate}{prvEx,prvAll}
\tagitem{prvEx}{%
  \iftag{probEx}{অনুশীলনী।}{\olref[syn][ass]{prop:sat-quant} অনুসারে
    $\Sat{M(\Gamma^*)}{\lexists[x][!A(x)]}$ যদি এবং কেবল যদি অন্তত একটি
    চল-মূল্যায়ন~$s$-এর জন্য $\Sat{M(\Gamma^*)}{!A(x)}[s]$। যেহেতু
    $\Domain{M(\Gamma^*)}$, $\Lang{L}$-এর বদ্ধ পদগুলি নিয়ে গঠিত, তাই
    এটি সত্য যদি এবং কেবল যদি অন্তত একটি বদ্ধ পদ~$t$ আছে যাতে $s(x) = t$
    এবং $\Sat{M(\Gamma^*)}{!A(x)}[s]$।
    \olref[fol][syn][ext]{prop:ext-formulas} অনুসারে, $s(x) = t$ হলে
    $\Sat{M(\Gamma^*)}{!A(x)}[s]$ যদি এবং কেবল যদি
    $\Sat{M(\Gamma^*)}{!A(t)}[s]$। \olref[fol][syn][ass]{prop:sentence-sat-true}
    অনুসারে $\Sat{M(\Gamma^*)}{!A(t)}[s]$ যদি এবং কেবল যদি
    $\Sat{M(\Gamma^*)}{!A(t)}$, কারণ $!A(t)$ একটি বাক্য।}}{}
\tagitem{prvAll}{%
  \iftag{probAll}{অনুশীলনী।}{\olref[syn][ass]{prop:sat-quant} অনুসারে
    $\Sat{M(\Gamma^*)}{\lforall[x][!A(x)]}$ যদি এবং কেবল যদি প্রতিটি
    চল-মূল্যায়ন~$s$-এর জন্য $\Sat{M(\Gamma^*)}{!A(x)}[s]$। মনে রাখি,
    $\Domain{M(\Gamma^*)}$, $\Lang{L}$-এর বদ্ধ পদগুলি নিয়ে গঠিত। তাই
    প্রতিটি বদ্ধ পদ~$t$-এর জন্য $s(x) = t$ এমন একটি চল-মূল্যায়ন, এবং
    যেকোনো চল-মূল্যায়নের ক্ষেত্রে $s(x)$ কোনো বদ্ধ পদ~$t$।
    \olref[fol][syn][ext]{prop:ext-formulas} অনুসারে, $s(x) = t$ হলে
    $\Sat{M(\Gamma^*)}{!A(x)}[s]$ যদি এবং কেবল যদি
    $\Sat{M(\Gamma^*)}{!A(t)}[s]$। \olref[fol][syn][ass]{prop:sentence-sat-true}
    অনুসারে $\Sat{M(\Gamma^*)}{!A(t)}[s]$ যদি এবং কেবল যদি
    $\Sat{M(\Gamma^*)}{!A(t)}$, কারণ $!A(t)$ একটি বাক্য।}}{}
\end{tagenumerate}
\end{proof}
}{}

\tagprob[FOL]{probEx,probAll}
\begin{prob}
\olref[fol][com][mod]{prop:quant-termmodel}-এর প্রমাণ সম্পূর্ণ করো।
\end{prob}
\tagendprob

\begin{lem}[সত্যতা-সহায়ক উপপাদ্য]
\ollabel{lem:truth} \iftag{FOL}{ধরা যাক $!A$-তে~$\eq$ নেই। তাহলে}{}
$\iftag{FOL}{\Sat{M(\Gamma^*)}}{\pSat{v(\Gamma^*)}}{!A}$ যদি এবং কেবল
যদি $!A \in \Gamma^*$।
\end{lem}

\begin{proof}
দুই অভিমুখ একই সঙ্গে, $!A$-এর ওপর আরোহের সাহায্যে প্রমাণ করি।
\begin{enumerate}
\tagitem{prvFalse}{\indcase{!A}{\lfalse}
  {পরিতৃপ্তির সংজ্ঞা অনুসারে
    $\iftag{FOL}{\Sat/{M(\Gamma^*)}{\lfalse}}{\pSat/{v(\Gamma^*)}{\lfalse}}$।
    অন্যদিকে, $\Gamma^*$ সঙ্গত বলে $\lfalse \notin \Gamma^*$।}}{}

\tagitem{prvTrue}{\indcase{!A}{\ltrue}
  {পরিতৃপ্তির সংজ্ঞা অনুসারে
    $\iftag{FOL}{\Sat{M(\Gamma^*)}}{\pSat{v(\Gamma^*)}}{\ltrue}$।
    অন্যদিকে, $\Gamma^*$ সঙ্গত ও !!{complete} এবং $\Gamma^* \Proves
    \ltrue$ বলে $\ltrue \in \Gamma^*$।}}{}

\tagitem{FOL}{\indcase{!A}{R(t_1, \dots, t_n)}
  {$\Sat{M(\Gamma^*)}{\Atom{R}{t_1, \dots, t_n}}$ যদি এবং কেবল যদি
    $\tuple{t_1, \dots, t_n} \in \Assign{R}{M(\Gamma^*)}$ (পরিতৃপ্তির
    সংজ্ঞা অনুসারে), যদি এবং কেবল যদি $R(t_1, \dots, t_n) \in
    \Gamma^*$ ($\Struct M(\Gamma^*)$-র নির্মাণ অনুসারে)।}}{\indcase{!A}{p}
    {$\pSat{v(\Gamma^*)}{p}$ যদি এবং কেবল যদি
    $\pAssign v(\Gamma^*)(p) = \True$ (পরিতৃপ্তির সংজ্ঞা অনুসারে), যদি
    এবং কেবল যদি $p \in \Gamma^*$ ($\pAssign v(\Gamma^*)$-র নির্মাণ
    অনুসারে)।}}

\tagitem{prvNot}{%
  \indcase{!A}{\lnot !B}
    {$\iftag{FOL}{\Sat{M(\Gamma^*)}}{\pSat{v(\Gamma^*)}}{\indfrm}$ যদি
    এবং কেবল যদি
    $\iftag{FOL}{\Sat/{M(\Gamma^*)}{!B}}{\pSat/{v(\Gamma^*)}{!B}}$
    (পরিতৃপ্তির সংজ্ঞা অনুসারে)। আরোহী পূর্বধারণা অনুসারে
    $\iftag{FOL}{\Sat/{M(\Gamma^*)}{!B}}{\pSat/{v(\Gamma^*)}{!B}}$ যদি
    এবং কেবল যদি $!B \notin \Gamma^*$। $\Gamma^*$ সঙ্গত ও !!{complete}
    বলে $!B \notin \Gamma^*$ যদি এবং কেবল যদি $\lnot !B \in
    \Gamma^*$।}}{}

\tagitem{prvAnd}{%
\iftag{probAnd}{%
  \indcase!{!A}{!B \land !C}{}}{%
  \indcase{!A}{!B \land
    !C}{$\iftag{FOL}{\Sat{M(\Gamma^*)}}{\pSat{v(\Gamma^*)}}{\indfrm}$
    যদি এবং কেবল যদি
    $\iftag{FOL}{\Sat{M(\Gamma^*)}}{\pSat{v(\Gamma^*)}}{!B}$ এবং
    $\iftag{FOL}{\Sat{M(\Gamma^*)}}{\pSat{v(\Gamma^*)}}{!C}$ উভয়ই
    সত্য (পরিতৃপ্তির সংজ্ঞা অনুসারে), যদি এবং কেবল যদি $!B \in
    \Gamma^*$ এবং $!C \in \Gamma^*$ উভয়ই সত্য (আরোহী পূর্বধারণা
    অনুসারে)। \olref[ccs]{prop:ccs}\olref[ccs]{prop:ccs-and} অনুসারে এটি
    সত্য যদি এবং কেবল যদি $(!B \land !C) \in \Gamma^*$।}}}{}

\tagitem{prvOr}{%
\iftag{probOr}{%
  \indcase!{!A}{!B \lor !C}{}}{%
  \indcase{!A}{!B \lor !C}
    {$\iftag{FOL}{\Sat{M(\Gamma^*)}}{\pSat{v(\Gamma^*)}}{\indfrm}$
    যদি এবং কেবল যদি
    $\iftag{FOL}{\Sat{M(\Gamma^*)}}{\pSat{v(\Gamma^*)}}{!B}$ অথবা
    $\iftag{FOL}{\Sat{M(\Gamma^*)}}{\pSat{v(\Gamma^*)}}{!C}$
    (পরিতৃপ্তির সংজ্ঞা অনুসারে), যদি এবং কেবল যদি $!B \in \Gamma^*$
    অথবা $!C \in \Gamma^*$ (আরোহী পূর্বধারণা অনুসারে)।
    \olref[ccs]{prop:ccs}\olref[ccs]{prop:ccs-or} অনুসারে এটি সত্য যদি এবং
    কেবল যদি $(!B \lor !C) \in \Gamma^*$।}}}{}

\tagitem{prvIf}{%
\iftag{probIf}{%
  \indcase!{!A}{!B \lif !C}{}}{%
  \indcase{!A}{!B \lif !C}
    {$\iftag{FOL}{\Sat{M(\Gamma^*)}}{\pSat{v(\Gamma^*)}}{\indfrm}$
    যদি এবং কেবল যদি
    $\iftag{FOL}{\Sat/{M(\Gamma^*)}}{\pSat/{v(\Gamma^*)}}{!B}$ অথবা
    $\iftag{FOL}{\Sat{M (\Gamma^*)}}{\pSat{v(\Gamma^*)}}{!C}$
    (পরিতৃপ্তির সংজ্ঞা অনুসারে), যদি এবং কেবল যদি $!B \notin
    \Gamma^*$ অথবা $!C \in \Gamma^*$ (আরোহী পূর্বধারণা অনুসারে)।
    \olref[ccs]{prop:ccs}\olref[ccs]{prop:ccs-if} অনুসারে এটি সত্য যদি
    এবং কেবল যদি $(!B \lif !C) \in \Gamma^*$।}}}{}

\iftag{FOL}{%
\tagitem{prvAll}{%
  \iftag{probAll}{%
    \indcase!{!A}{\lforall[x][!B(x)]}{}}{%
    \indcase{!A}{\lforall[x][!B(x)]}{$\Sat{M(\Gamma^*)}{\indfrm}$ যদি
      এবং কেবল যদি সব বদ্ধ পদ~$t$-এর জন্য
      $\Sat{M(\Gamma^*)}{!B(t)}$ (\olref{prop:quant-termmodel})। আরোহী
      পূর্বধারণা অনুসারে এটি সত্য যদি এবং কেবল যদি সব বদ্ধ পদ~$t$-এর
      জন্য $!B(t) \in \Gamma^*$; আর \olref[hen]{prop:saturated-instances}
      অনুসারে এটি সত্য যদি এবং কেবল যদি $\lforall[x][!B(x)] \in
      \Gamma^*$।}}}{}

\tagitem{prvEx}{%
  \iftag{probEx}{%
    \indcase!{!A}{\lexists[x][!B(x)]}{}}{%
    \indcase{!A}{\lexists[x][!B(x)]}{$\Sat{M(\Gamma^*)}{\indfrm}$ যদি
      এবং কেবল যদি অন্তত একটি বদ্ধ পদ~$t$-এর জন্য
      $\Sat{M(\Gamma^*)}{!B(t)}$ (\olref{prop:quant-termmodel})। আরোহী
      পূর্বধারণা অনুসারে এটি সত্য যদি এবং কেবল যদি অন্তত একটি বদ্ধ
      পদ~$t$-এর জন্য $!B(t) \in \Gamma^*$।
      \olref[hen]{prop:saturated-instances} অনুসারে এটি সত্য যদি এবং কেবল
      যদি $\lexists[x][!B(x)] \in \Gamma^*$।}}}{}% BN-SRC-077 TARGET-CORRECTION-BEGIN
\par\smallskip\noindent\textnormal{\small\textbf{উৎস-স্পষ্টীকরণ (BN-SRC-077):} হিমায়িত ইংরেজি সত্যতা-সহায়ক উপপাদ্যের পরিমাণসূচক-ক্ষেত্রগুলি “সব পদ” ও “কোনো পদ” বলে, কিন্তু যে \olref{prop:quant-termmodel} এবং \olref[hen]{prop:saturated-instances} ব্যবহার করা হয়েছে উভয়ই বদ্ধ পদের ওপর পরিমাণায়িত। বাংলা পাঠে দুই ক্ষেত্রেই সেই বদ্ধ-পদ পরিসর রাখা হয়েছে।} % BN-SRC-077 TARGET-CORRECTION-END% BN-SRC-078 TARGET-CORRECTION-BEGIN
\par\smallskip\noindent\textnormal{\small\textbf{উৎস-সংশোধন (BN-SRC-078):} হিমায়িত ইংরেজি সার্বিক-পরিমাণসূচক ক্ষেত্রে আরোহী সূত্রকে $\lforall[x][!B(x)]$ ধরেও শেষ সদস্যতা-দাবিতে ভুল করে $\lforall[x][!A(x)]$ লেখে। বাংলা লক্ষ্যে সর্বত্র সংশ্লিষ্ট উপসূত্র~$!B(x)$ রাখা হয়েছে।} % BN-SRC-078 TARGET-CORRECTION-END
}{}
\end{enumerate}
\end{proof}


\tagprob[FOL]{probOr,probAnd,probIf,probEx,probAll}

\begin{prob}
\olref[fol][com][mod]{lem:truth}-এর প্রমাণ সম্পূর্ণ করো।
\end{prob}

\tagendprob

\tagprob[notFOL]{probOr,probAnd,probIf}

\begin{prob}
\olref[pl][com][mod]{lem:truth}-এর প্রমাণ সম্পূর্ণ করো।
\end{prob}

\tagendprob

\end{document}
